\Needspace{18\baselineskip}
\section{Independent deep substitution and pointed host classes}
\label{sec:deep-interchange}
This section completes the host/site quotient and independent-site
interchange obligation formerly listed in Appendix~\ref{app:remaining-work}.
The theorem constructs surviving addresses after each substitution, at
arbitrary depth, and compares the final graphs by recursive isomorphism.
It complements nested associativity: independence excludes replacing a
box and one of its descendants as the two independent steps.
No hole-freshness or intermediate context-validity assumption is needed
for the graph isomorphism.

\Needspace{8\baselineskip}
\subsection{Interchange at distinct shallow edges}
Let $e\ne f$ be edges of a finite typed graph $C$, and let $H$ and $K$
have their respective interfaces. After replacing $e$, write $f/e$ for
the retained occurrence of $f$, and define $e/f$ symmetrically.
\begin{theorem}[Shallow interchange]
There is an explicit ordered-interface-preserving isomorphism
\[
       (C[e:=H])[f/e:=K]\ \cong\ (C[f:=K])[e/f:=H].
\]
The host may contain cycles or self-loops.
\end{theorem}
\begin{proof}
The internal carriers are $(V_C+V_H)+V_K$ and $(V_C+V_K)+V_H$.
The bijection fixes the original host component, sends the first filling
component to its new position, and sends the second filling component to
its new position. The inverse performs the reverse rearrangement;
both identities follow by the three cases.
The edge bijection similarly fixes the surviving original edges, each
different from both selected edges, and exchanges the positions of the
two filling-edge carriers. Its subtype proofs follow from $e\ne f$ and
the fact that the original surviving edge differs from each selected edge.
Again the three cases give the inverse identities.
\lean{Graph/PlugInterchange}{12}{17}{Graph.Net.PlugInterchange.vertices}
\lean{Graph/PlugInterchange}{19}{32}{Graph.Net.PlugInterchange.edges}

For any retained edge, its source and target attachment ports after one
replacement are the original ports mapped through host inclusion.
Taking an indexed element of the mapped incidence list proves these
two statements.
\lean{Graph/PlugInterchange}{34}{36}{Graph.Net.PlugInterchange.sourcePort_retained}
\lean{Graph/PlugInterchange}{38}{40}{Graph.Net.PlugInterchange.targetPort_retained}

The composite inclusion of a host vertex commutes with the carrier
bijection, by cases on internal, input, and output vertices.
For a vertex of $H$, its inclusion after inserting $H$ first agrees with
its inclusion when $H$ is inserted second. The internal case is
definitional. The input and output cases use the retained-port formula
for $e$ after inserting $K$, followed by the host-vertex identity.
For $K$ the argument is reversed: first use its retained-port formula
after inserting $H$, then the host identity. These three commuting maps
cover every vertex used by an incidence list.
\lean{Graph/PlugInterchange}{42}{46}{Graph.Net.PlugInterchange.map_host}
\lean{Graph/PlugInterchange}{48}{63}{Graph.Net.PlugInterchange.map_first}
\lean{Graph/PlugInterchange}{65}{79}{Graph.Net.PlugInterchange.map_second}

Labels are unchanged under the carrier rearrangement. For source and
target lists, split an edge into surviving host, $H$, or $K$ origin;
compose the list maps and apply the corresponding commuting identity
pointwise. This verifies ordered incidence and completes the
isomorphism, not merely a bijection of carriers.
\end{proof}
\lean{Graph/PlugInterchange}{82}{116}{Graph.Net.PlugInterchange.iso}

\Needspace{8\baselineskip}
\subsection{Changing a box beside another substitution}
A change to the body of box edge $f$ leaves the label of every $e\ne f$
unchanged. Unfolding the conditional label update proves this directly.
The source and target interfaces at $e$ are therefore equal to their
original interfaces. The surviving shallow site is obtained by transporting
its interface indices along these equalities.
\lean{Graph/BoxInterchange}{9}{14}{Linear.replaceBox_other_label}
\lean{Graph/BoxInterchange}{17}{26}{Linear.Site.afterBox}

\begin{lemma}[Mixed shallow and box interchange]
Replacing a shallow edge $e$ and changing the body of a distinct box edge
$f$ commute up to ordinary graph isomorphism, with the shallow site
transported to its unchanged interface.
\end{lemma}
\begin{proof}
A body update preserves the host's entire outer incidence structure.
To reuse the general incidence-transport construction, first regard its
edge labels as related by the always-true relation. Identity carrier
bijections then give an incidence isomorphism from the original host to
the updated host. This auxiliary relation is used only for transporting
incidence; it is not the final equality condition on labels.
\lean{Graph/BoxInterchange}{29}{38}{Linear.replaceBox_incidence}

Lift that incidence isomorphism through substitution at $e$, transporting
the replacement's boundary indices along the equal interfaces.
Use its vertex and edge bijections, internal labels, and source and target
proofs for the desired final isomorphism. Check actual edge-label equality
separately: a retained original edge has its updated box label precisely
when it is $f$, while every inserted edge has its original replacement
label. Equality with $f$ is equivalent to equality of its retained copy
with the retained copy of $f$. This proves literal label equality and
hence the stated ordinary graph isomorphism.
\end{proof}
\lean{Graph/BoxInterchange}{41}{73}{Linear.plug_retained_box_iso}

\begin{lemma}[Distinct box updates commute]
Changing the bodies at two distinct box edges, in either order, yields
equal graph records.
\end{lemma}
\begin{proof}
All carriers and incidence fields are unchanged. For the label function,
split an edge into the cases that it is the first selected edge, the
second selected edge, or neither. Distinctness excludes the two updates
interfering, so the same final label is obtained in every case.
\end{proof}
\lean{Graph/BoxInterchange}{76}{93}{Linear.replaceBox_comm}

\Needspace{8\baselineskip}
\subsection{Structural independence and residual addresses}
\begin{definition}[Independent sites]\label{term:independent-sites}
A certificate of independence is generated structurally as follows.
Two distinct shallow edges are independent. A shallow edge and a site
inside a different box are independent. Sites inside two different
boxes are independent, regardless of their inner depths. Two sites inside
the same box are independent when their inner sites are independent.
Finally independence can be reversed. These constructors allow any finite
common chain of enclosing boxes before the addresses diverge; they give
no constructor relating a selected edge to its own descendant.
The certificate contains structural occurrence information, not an
assumed interchange theorem.
\end{definition}
\lean{Graph/IndependentSites}{10}{30}{Linear.IndependentSites}

\begin{definition}[Interchange result]\label{term:interchange-result}
For independent $s,t$ in $C$ and replacements $A,B$, an interchange result
contains a residual site $t/s$ in $C[s:=A]$, a residual site $s/t$ in
$C[t:=B]$, and a recursive isomorphism
\[
        (C[s:=A])[t/s:=B]\ \simeq\
        (C[t:=B])[s/t:=A].
\]
Here the residual sites are constructed by following retained occurrences,
including their enclosing boxes. They are not chosen by searching for
graphs with equal results.
\end{definition}
\lean{Graph/IndependentSites}{34}{39}{Linear.InterchangeResult}

\begin{theorem}[Independent deep interchange]
Every structural independence certificate and pair of interface-compatible
replacements produces an interchange result.
\end{theorem}
\begin{proof}
Induct on the independence certificate, with replacements arbitrary.
For distinct shallow edges, use their retained copies and the shallow
interchange isomorphism. For a shallow edge and a site inside another
box, retain that box address after the shallow plug; in the other order,
use the shallow site transported across the body update. The mixed
interchange lemma supplies the isomorphism.

For sites in different boxes, retain each enclosing edge and inner
address after updating the other box. The unchanged-label lemma supplies
the necessary box-label certificates. Distinct box updates commute as
records, so recursive reflexivity gives the comparison.

For two sites in a shared box, apply the induction hypothesis in its
body. Lift each residual address through the updated enclosing box.
Performing the second substitution updates that same box a second time;
the overwrite lemma reduces each side to a single update with the final
inner body. Lift the inner recursive isomorphism through this common
outer update. This case is what permits arbitrary shared nesting depth.

For a reversed independence certificate, apply the induction hypothesis
with replacements reversed, exchange its two residual sites, and invert
its isomorphism. All cases construct both addresses and their comparison.
\end{proof}
\lean{Graph/IndependentSites}{43}{79}{Linear.IndependentSites.interchange}

Transporting a site's interface along equalities does not change its
selected hole label: eliminate the equalities and the claim is reflexive.
Induction on independence then proves that both residual addresses
select the original hole labels. The shallow and separate-box cases
are definitional; the mixed case uses interface transport and the
unchanged-label lemma; the shared-box case uses the induction
hypothesis; reversing independence exchanges the two assertions.
This also covers non-hole sites, whose selected label is absent.
\lean{Graph/IndependentSites}{82}{85}{Linear.Site.selectedHole_reinterface}
\lean{Graph/IndependentSites}{88}{104}{Linear.IndependentSites.residual_selectedHole}

The two final recursive graph classes are equal by quotient soundness
applied to the constructed isomorphism. They have unique holes
simultaneously, by invariance of that property under recursive
isomorphism. The latter equivalence concerns the final graphs; neither
intermediate graph is required to be a valid context.
\lean{Graph/IndependentSites}{107}{112}{Linear.IndependentSites.interchange_modulo}
\lean{Graph/IndependentSites}{115}{120}{Linear.IndependentSites.unique_iff}

\Needspace{8\baselineskip}
\subsection{A quotient of hosts with selected occurrences}
\begin{definition}[Pointed hypernet]\label{term:pointed-hypernet}
A pointed hypernet consists of a host $C:I\to O$ at depth bound $d$ and
a selected site $s$ with local depth $r$ and interface $J\to K$.
All these type indices are fixed for a particular space of pointed
hypernets, while the host and its carrier names may vary.
\end{definition}
\lean{Graph/SiteQuotient}{9}{12}{Linear.PointedHypernet}

\begin{definition}[Pointed renaming]\label{term:pointed-renaming}
Two pointed hypernets are equivalent under the equivalence relation
generated by structural site correspondences. Its constructors are a
single existing $\mathsf{SiteIso}$ correspondence, reflexivity, reversal,
and composition. Thus the relation consists of finite renaming zigzags.
Equality of substituted results is not a generator: the selected occurrence
remains part of the object being quotiented.
This construction does not assert that the previously defined
$\mathsf{SiteIso}$ itself has already been proved to be an equivalence.
\end{definition}
\lean{Graph/SiteQuotient}{16}{21}{Linear.PointedRenaming}

A single site correspondence includes an outer recursive host isomorphism:
inspect either its shallow or inside constructor. A finite pointed
renaming therefore relates recursively isomorphic hosts, by induction;
identity, inverse, and composition of graph isomorphisms handle the
three equivalence constructors.
\lean{Graph/SiteQuotient}{24}{29}{Linear.SiteIso.host}
\lean{Graph/SiteQuotient}{31}{38}{Linear.PointedRenaming.host}

Pointed renaming preserves plugging of a fixed replacement. For one
structural correspondence use the existing site-substitution theorem.
For reflexivity, reversal, and composition use the corresponding
recursive isomorphism operations on results. If the replacements are
also recursively isomorphic, compose this result with replacement
congruence at the second selected site.
\lean{Graph/SiteQuotient}{41}{49}{Linear.PointedRenaming.plug}
\lean{Graph/SiteQuotient}{51}{56}{Linear.PointedRenaming.plug_congr}

Retained-hole inventories also agree along any pointed renaming.
The single-step case is the existing site correspondence invariant;
the other cases use reflexivity, symmetry, and transitivity of logical
equivalence.
\lean{Graph/SiteQuotient}{59}{66}{Linear.PointedRenaming.retainedHole_iff}

\begin{definition}[Pointed isomorphism class]\label{term:pointed-class}
Use pointed renaming as a setoid relation, with its three equivalence
constructors supplying the equivalence laws. A pointed isomorphism
class is its quotient. The quotient identifies hosts and their selected
occurrences together, with outer and local interfaces fixed.
\end{definition}
\lean{Graph/SiteQuotient}{68}{71}{Linear.pointedSetoid}
\lean{Graph/SiteQuotient}{74}{75}{Linear.PointedModuloIso}

\begin{definition}[Fully quotiented substitution]\label{term:pointed-plug}
Given a pointed class $[C,s]$ and a replacement class $[A]$, define
\[
          \mathsf{pointedPlug}([C,s],[A])=[C[s:=A]].
\]
For a fixed pointed representative, the previously constructed
replacement-quotient operation already makes the right-hand side
independent of the representative of $[A]$. To eliminate the pointed
quotient, use function extensionality and quotient induction on $[A]$.
Pointed-renaming substitution invariance then proves equality of the
result classes. On concrete representatives the defining equation is
reflexive.
\end{definition}
\lean{Graph/SiteQuotient}{78}{85}{Linear.pointedPlug}
\lean{Graph/SiteQuotient}{87}{89}{Linear.pointedPlug_mk}

Forgetting the selected site gives a well-defined host class, because
pointed renaming preserves recursive host isomorphism.
The retained-hole predicate likewise descends to pointed classes,
using the retained-inventory equivalence and propositional extensionality.
\lean{Graph/SiteQuotient}{92}{94}{Linear.pointedHost}
\lean{Graph/SiteQuotient}{97}{99}{Linear.pointedRetainedHole}

\begin{definition}[Quotient freshness]\label{term:pointed-freshness}
The predicate $\mathsf{pointedFresh}([C,s],[A])$ means that no hole
retained outside $s$ occurs in $A$. Replacement-renaming invariance
permits eliminating the replacement quotient. Retained-inventory
invariance permits eliminating the pointed quotient: in either
direction, transfer the retained-hole witness and apply the original
freshness condition. Function and propositional extensionality turn
these equivalences into the equalities required by quotient elimination.
\end{definition}
\lean{Graph/SiteQuotient}{102}{112}{Linear.pointedFresh}

\begin{theorem}[Sharp context validity on classes]
If the host class of $P$ has unique holes, then
\[
 \mathsf{UniqueHoles}(\mathsf{pointedPlug}(P,[A]))
 \quad\Longleftrightarrow\quad
 \mathsf{UniqueHoles}([A])\ \land\ \mathsf{pointedFresh}(P,[A]).
\]
\end{theorem}
\begin{proof}
Apply quotient induction to both arguments. The assertion becomes the
existing sharp context criterion at the concrete selected site, with
the supplied host-uniqueness hypothesis. No stronger disjointness
condition on the whole host is introduced.
\end{proof}
\lean{Graph/SiteQuotient}{115}{122}{Linear.pointedPlug_unique_iff}

Transporting a site along a recursive host isomorphism leaves its pointed
class unchanged: the transport specification is a single structural
site correspondence, hence a generator of pointed renaming.
Finally, applying fully quotiented substitution at either of the
constructed residual pointed hosts gives equal result classes by
deep interchange and quotient soundness. This is the independent
substitution law expressed through the new operation.
\lean{Graph/SiteQuotient}{125}{129}{Linear.pointed_transport}
\lean{Graph/SiteQuotient}{132}{137}{Linear.IndependentSites.pointed_interchange}

\Needspace{8\baselineskip}
\subsection{Examples that distinguish occurrences}
The second site in the earlier boxed two-hole chain selects its second
inner hole. The two sites are independent by the common-box constructor
applied to distinct shallow holes. Filling one by the repeated-hole
fragment and the other by a variable fragment gives equal final
classes in either order. The earlier collision theorem proves that one
intermediate graph has a duplicated hole label; interchange still
applies because it is a theorem about well-formed graphs, not only valid
intermediate contexts.
\lean{Graph/InterchangeExamples}{9}{10}{Linear.Examples.boxedSecondSite}
\lean{Graph/InterchangeExamples}{12}{13}{Linear.Examples.boxedSitesIndependent}
\lean{Graph/InterchangeExamples}{17}{22}{Linear.Examples.boxed_interchange}

A second host has two outer thunk inputs, no outputs or internal
vertices, and two distinct box edges. Both edges have exactly the same
single-hole body. Each source list contains its own input; both target
lists are empty. Injectivity of the input constructor proves source
uniqueness, while target uniqueness is vacuous. Inputs witness source
coverage; there are no other vertices requiring coverage. Splitting
the two-element edge type proves the source typings, and the empty
target typings are reflexive. This constructs an actual finite graph.
It need not be a unique-hole context, since the bodies deliberately
repeat the same hole label.
\lean{Graph/InterchangeExamples}{25}{59}{Linear.Examples.twinBoxes}

Select the sole inner hole of either box. The separate-box constructor
certifies independence despite equality of the box labels. Fill the
first with the value operation and the second with the other-value
operation. Deep interchange proves agreement of the two orders; each
residual address follows its own enclosing edge.
\lean{Graph/InterchangeExamples}{61}{62}{Linear.Examples.twinSite}
\lean{Graph/InterchangeExamples}{64}{66}{Linear.Examples.twinIndependent}
\lean{Graph/InterchangeExamples}{69}{72}{Linear.Examples.twin_interchange}

The correctly transported first site in the renamed boxed chain gives
the same pointed class as the original site, by the previously checked
structural correspondence. In contrast, keeping its former numerical
index selects the wrong hole and gives a different pointed class.
If those classes were equal, quotient exactness would give a pointed
renaming. Substitution invariance with the repeated-hole replacement
would then give recursively isomorphic results, contradicting the
existing wrong-address counterexample.
\lean{Graph/InterchangeExamples}{75}{80}{Linear.Examples.pointed_renamed_first}
\lean{Graph/InterchangeExamples}{83}{89}{Linear.Examples.pointed_wrong_address}

The proof dependency is explicit: shallow carrier interchange and local
box-update laws yield arbitrary-depth residual construction; site
correspondence invariants yield the pointed quotient and its substitution
operation; recursive isomorphism connects interchange to equality of
result classes. Deep specimen representations and the concrete UAM
rules are still separate obligations in Appendix~\ref{app:remaining-work}.
